% --- Back cover ---
% Large PERSONIX wordmark, no logo image. Cream paper inherited globally.
% Absolute (tikz overlay) positioning so the footer sticks to the bottom.
\clearpage
\thispagestyle{empty}
\begin{tikzpicture}[remember picture, overlay]
  % --- Wordmark ---
  \node[anchor=north, font=\sffamily\bfseries\fontsize{48}{52}\selectfont]
    at ([yshift=-26mm]current page.north) {PERSONIX};
  \node[anchor=north, font=\rmfamily\itshape\large]
    at ([yshift=-44mm]current page.north) {Ödünsüz Değişim};
  \node[anchor=north, font=\rmfamily\small,
        text width=\paperwidth-50mm, align=center]
    at ([yshift=-54mm]current page.north)
    {Sansürlenemez ve Yozlaştırılamaz Merkeziyetsiz Bir İtibar Ağı};
  % --- Body ---
  \node[anchor=north, text width=\paperwidth-50mm, align=justify, font=\small]
    at ([yshift=-72mm]current page.north)
    {%
      Devlet, iletişim yavaşken, kararlar yerelken ve otorite yönettiği
      kasabada yaşarken işe yaradı. Bugünün ağları bu üç varsayımı da tersine
      çeviriyor --- ve üzerlerine kurulan kurumlar artık yönettikleri dünyaya
      uymuyor. Bu kitap uyan bir aracı anlatıyor: kimliğin senin olduğu,
      gerçeğin kamuya açık biçimde denetlenebildiği ve rüşvetin itibar
      açısından intihar olduğu merkeziyetsiz bir itibar ağı.

      \smallskip
      Bir manifesto değil. Bir öngörü değil. Devletin ardılının nasıl inşa
      edilebileceğine dair işleyen bir plan --- gönüllü olarak, her seferinde
      bir bildirimle.
    };
  % --- Footer (anchored to the bottom of the page) ---
  \node[anchor=south, font=\sffamily\footnotesize,
        text width=\paperwidth-50mm, align=center]
    at ([yshift=12mm]current page.south)
    {%
      Pavel Kudrna\quad\textbullet\quad Prag, 2026\par
      \href{https://personix.org}{personix.org}\par
      Creative Commons Attribution-ShareAlike 4.0 International
      (CC BY-SA 4.0) lisansı altındadır.
    };
\end{tikzpicture}%
\mbox{}%
\clearpage
